\documentclass[a5paper,10pt]{article}

% https://github.com/InDevRus/filippov-solutions

% Нумерация страниц
\pagenumbering{gobble}

% Настройка полей
\usepackage{geometry}
\geometry{left=1cm}
\geometry{right=1cm}
\geometry{top=1cm}
\geometry{bottom=1.5cm}

% Вставка таблицы
\usepackage{array}
\usepackage{tabularx}

% Инструменты выравнивания формул
\usepackage{amsmath}
\usepackage{mathtools}

% Диакритика в математике
\usepackage{amsfonts}

% Вычисления размеров на ходу
\usepackage{calc}

% Удобные записи производных
\usepackage[italic=true]{derivative}

% Настройки сносок
\usepackage{enotez}

% Длинные знаки векторов
\usepackage{anyfontsize}
\usepackage[b]{esvect}

% Вставка картинок
\usepackage{graphicx}

% Вставка красной строки
\usepackage{indentfirst}
\setlength{\parindent}{1em}

% Условия в командах
\usepackage{xifthen}

% Луа-скрипты
\usepackage{luacode}

% Настройка команд отдельно в математическом и текстовом режимах
\usepackage{mathcommand}

% Для повторения LaTeX кода
\usepackage{multido}

% Конфигурация отступов формул
\delimitershortfall-1sp
\usepackage{mleftright}
\mleftright

% Растягивание объектов
\usepackage{relsize}

% Обтекание текстом
\usepackage{wrapfig}
\usepackage{subcaption}
\usepackage{floatflt}

% Поддержка ввода кириллицы
\usepackage[english, russian]{babel}

% Настройка корневой позиции для компилятора
\usepackage{import}
\providecommand*{\ProjectRootPath}{..}

\import{\ProjectRootPath/preambles/macros/}{theorems}

\import{\ProjectRootPath/preambles/fonts}{font_setup}

% Знак градуса
\usepackage{siunitx}

\import{\ProjectRootPath/preambles/macros/}{operators}
\import{\ProjectRootPath/preambles/macros/}{redefinitions}

\import{\ProjectRootPath/preambles/fonts}{letters_setup}
\import{\ProjectRootPath/preambles/fonts/adjustments}{custom_superscripts}

\import{\ProjectRootPath/preambles/custom}{formatting}
\import{\ProjectRootPath/preambles/custom}{frames}
\import{\ProjectRootPath/preambles/custom}{list_setup}

% TODO: перевести команды на современный стиль объявления

\begin{document}

\begin{align*} 
    \der{}{t} \begin{pmatrix} \,x\br{t} \\ \,y\br{t} \end{pmatrix} 
        = A \begin{pmatrix} \,x\br{t} \\ \,y\br{t} \end{pmatrix} 
        + \begin{pmatrix} o\br{\vbr{X}} \\ o\br{\vbr{X}} \end{pmatrix}
    && A = \begin{pmatrix} a & 1 \\ 1 & a \end{pmatrix}
\end{align*}

$\det\br{A - \lambda E} = \br{\lambda - a}^2 - 1 = 0$.
$D > 0$ для всех $a$, так что корни веществены и различны.
$\pow{\lambda}{2} - 2a \lambda + \pow{a}{2} - 1 = 0$.
\begin{itemize}
    \item при $a \in \br{-1; 1}$ будет
        $\sub{\lambda}{1} \cdot \sub{\lambda}{2} = \pow{a}{2} - 1 < 0$,
        то есть собственные значения разных знаков, и нулевое решение неустойчиво;
    
    \item при $a \in \br{-\infty; -1}$ имеем 
        $\sub{\lambda}{1} \cdot \sub{\lambda}{2} > 0$
        и $\sub{\lambda}{1} + \sub{\lambda}{2} = 2a < 0$,
        то есть оба собственных значения отрицательны,
        а нулевое решение устойчиво асимптотически;
    
    \item при $a \in \left[1; +\infty \right)$ получаем 
        $\sub{\lambda}{1} \cdot \sub{\lambda}{2} \ge 0$ 
        и $\sub{\lambda}{1} + \sub{\lambda}{2} > 0$,
        так что по меньшей мере одно из собственных значений положительно,
        и нулевое решение неустойчиво;
    
    \item при $a = -1$ собственные значения равны 
        $\sub{\lambda}{1} = 0$, $\sub{\lambda}{2} = -2$,
        и по первому приближению утверждать об устойчивости нулевого решения нельзя.
\end{itemize}

Для случая $a = -1$ покажем, что нулевое решение неустойчиво с помощью функции Ляпунова.
Рассмотрим функцию $v\br{x; y} = \pow{e}{x + y} - 1$ 
в открытой области $V = \cbr{x + y > 0}$. 
На границе $\partial V = \cbr{x + y = 0}$ имеем $v\br{x; y} = 0$,
внутри области $v\br{x; y} > 0$, но 
$$\der{v}{t}\bigg|_{a = -1}
    = \pow{e}{x + y} \br{\pow{x}{2} + y - x} + \pow{e}{x + y} \br{\pow{y}{2} - y + x} 
    = \pow{e}{x + y} \br{\pow{x}{2} + \pow{y}{2}} > 0,$$
так как точка $\br{0; 0}$ не лежит в $V$. По теореме Четаева нулевое решение неустойчиво.

Итак, нулевое решение устойчиво асимптотически только при $a < -1$.

\end{document}
